t%%%%%%%%%%%%%%%%%%%%%%%%%%%%%%%%%%%%%%%%%%%%%%%%%%%%%%%%%%%%
% 10.4 CONVEX OPTIMIZATION
% The Composition of Advanced Mathematics
%%%%%%%%%%%%%%%%%%%%%%%%%%%%%%%%%%%%%%%%%%%%%%%%%%%%%%%%%%%%

\section{Convex Optimization}
\label{sec:convex-optimization}

\prerequisite{
Optimization fundamentals, linear programming, nonlinear programming,
convex sets, convex functions, gradients, Hessians, affine geometry,
Lagrange multipliers, KKT conditions, and basic linear algebra.
}

\learningobjectives{
By the end of this section, the reader should be able to:
\begin{itemize}
    \item define convex sets and convex optimization problems;
    \item verify convexity using geometric and algebraic criteria;
    \item recognize affine sets, half-spaces, balls, cones, and
          intersections as convex sets;
    \item define convex combinations and convex hulls;
    \item define convex functions and strictly convex functions;
    \item use first-order and second-order convexity criteria;
    \item distinguish convexity from strict and strong convexity;
    \item define epigraphs and use them to characterize convexity;
    \item understand preservation rules for convex functions;
    \item recognize convexity under affine transformations;
    \item understand pointwise maxima and supremum constructions;
    \item recognize perspective functions;
    \item define convex cones;
    \item understand dual cones;
    \item recognize second-order and positive-semidefinite cones;
    \item formulate convex quadratic programs;
    \item understand conic optimization conceptually;
    \item recognize second-order cone programming;
    \item recognize semidefinite programming;
    \item understand why local minima of convex problems are global;
    \item understand uniqueness under strict convexity;
    \item formulate the Lagrange dual function;
    \item understand weak duality;
    \item define the dual problem;
    \item understand strong duality under Slater's condition;
    \item interpret KKT conditions as sufficient for convex problems;
    \item understand complementary slackness;
    \item understand supporting hyperplanes and subgradients;
    \item define subgradients for nonsmooth convex functions;
    \item understand subgradient optimality conditions;
    \item recognize proximal operators conceptually;
    \item understand projection onto convex sets;
    \item recognize projected gradient methods;
    \item understand gradient descent for smooth convex functions;
    \item understand strong convexity and convergence rates conceptually;
    \item recognize accelerated first-order methods conceptually;
    \item understand interior-point methods for convex programs;
    \item recognize barrier functions and self-concordance conceptually;
    \item understand convex relaxations;
    \item connect convex optimization with statistics, machine learning,
          signal processing, control, finance, and operations research.
\end{itemize}
}

%%%%%%%%%%%%%%%%%%%%%%%%%%%%%%%%%%%%%%%%%%%%%%%%%%%%%%%%%%%%
\subsection{Why Convex Optimization Is Special}
\label{subsec:why-convex-optimization-special}
%%%%%%%%%%%%%%%%%%%%%%%%%%%%%%%%%%%%%%%%%%%%%%%%%%%%%%%%%%%%

Convex optimization is the study of minimizing a convex function over a
convex feasible set.

A standard convex problem has form

\[
\boxed{
\begin{aligned}
\min_{\mathbf{x}}
\quad&
f_0(\mathbf{x})
\\
\text{subject to}
\quad&
f_i(\mathbf{x})\leq0,
\qquad
i=1,\ldots,m,
\\
&
A\mathbf{x}=\mathbf{b},
\end{aligned}
}
\]

where:

\[
f_0,f_1,\ldots,f_m
\]

are convex functions and the equality constraints are affine.

Convex optimization is especially important because

\[
\boxed{
\text{every local minimum is global}.
}
\]

%%%%%%%%%%%%%%%%%%%%%%%%%%%%%%%%%%%%%%%%%%%%%%%%%%%%%%%%%%%%
\subsection{Convex Combinations}
\label{subsec:convex-combinations}
%%%%%%%%%%%%%%%%%%%%%%%%%%%%%%%%%%%%%%%%%%%%%%%%%%%%%%%%%%%%

A convex combination of points

\[
\mathbf{x}_1,\ldots,\mathbf{x}_k
\]

is a vector of the form

\[
\boxed{
\sum_{i=1}^{k}
\theta_i\mathbf{x}_i,
}
\]

where

\[
\boxed{
\theta_i\geq0,
}
\]

and

\[
\boxed{
\sum_{i=1}^{k}
\theta_i=1.
}
\]

For two points,

\[
\boxed{
\theta\mathbf{x}
+
(1-\theta)\mathbf{y},
\qquad
0\leq\theta\leq1,
}
\]

describes the line segment between them.

%%%%%%%%%%%%%%%%%%%%%%%%%%%%%%%%%%%%%%%%%%%%%%%%%%%%%%%%%%%%
\subsection{Convex Sets}
\label{subsec:convex-sets-convex-optimization}
%%%%%%%%%%%%%%%%%%%%%%%%%%%%%%%%%%%%%%%%%%%%%%%%%%%%%%%%%%%%

\begin{definitionbox}{Convex Set}
A set \(C\subseteq\mathbb{R}^n\) is convex if every line segment between
two points of \(C\) lies entirely in \(C\).
\end{definitionbox}

Thus, for all

\[
\mathbf{x},\mathbf{y}\in C
\]

and

\[
0\leq\theta\leq1,
\]

we require

\[
\boxed{
\theta\mathbf{x}
+
(1-\theta)\mathbf{y}
\in C.
}
\]

%%%%%%%%%%%%%%%%%%%%%%%%%%%%%%%%%%%%%%%%%%%%%%%%%%%%%%%%%%%%
\subsection{Worked Example: A Convex Set}
\label{subsec:worked-convex-set}
%%%%%%%%%%%%%%%%%%%%%%%%%%%%%%%%%%%%%%%%%%%%%%%%%%%%%%%%%%%%

\begin{workedexamplebox}{Half-Space}

Consider

\[
C
=
\{
\mathbf{x}\in\mathbb{R}^n:
\mathbf{a}^T\mathbf{x}\leq b
\}.
\]

Take

\[
\mathbf{x},\mathbf{y}\in C.
\]

Then

\[
\mathbf{a}^T\mathbf{x}\leq b,
\]

and

\[
\mathbf{a}^T\mathbf{y}\leq b.
\]

For

\[
0\leq\theta\leq1,
\]

\[
\mathbf{a}^T
[
\theta\mathbf{x}
+
(1-\theta)\mathbf{y}
]
=
\theta\mathbf{a}^T\mathbf{x}
+
(1-\theta)\mathbf{a}^T\mathbf{y}.
\]

Therefore,

\[
\mathbf{a}^T
[
\theta\mathbf{x}
+
(1-\theta)\mathbf{y}
]
\leq
\theta b+(1-\theta)b
=
b.
\]

Hence,

\begin{answerbox}
\[
\boxed{
C
\text{ is convex}.
}
\]
\end{answerbox}

\end{workedexamplebox}

%%%%%%%%%%%%%%%%%%%%%%%%%%%%%%%%%%%%%%%%%%%%%%%%%%%%%%%%%%%%
\subsection{Affine Sets}
\label{subsec:affine-sets-convex}
%%%%%%%%%%%%%%%%%%%%%%%%%%%%%%%%%%%%%%%%%%%%%%%%%%%%%%%%%%%%

An affine set has form

\[
\boxed{
C
=
\{
\mathbf{x}:
A\mathbf{x}=\mathbf{b}
\}.
}
\]

Every affine set is convex.

Indeed, if

\[
A\mathbf{x}=\mathbf{b},
\qquad
A\mathbf{y}=\mathbf{b},
\]

then

\[
A
[
\theta\mathbf{x}
+
(1-\theta)\mathbf{y}
]
=
\mathbf{b}.
\]

%%%%%%%%%%%%%%%%%%%%%%%%%%%%%%%%%%%%%%%%%%%%%%%%%%%%%%%%%%%%
\subsection{Balls and Norm Constraints}
\label{subsec:balls-convex}
%%%%%%%%%%%%%%%%%%%%%%%%%%%%%%%%%%%%%%%%%%%%%%%%%%%%%%%%%%%%

For any norm \(\|\cdot\|\), the ball

\[
\boxed{
B
=
\{
\mathbf{x}:
\|\mathbf{x}-\mathbf{x}_0\|\leq r
\}
}
\]

is convex.

This follows from the triangle inequality:

\[
\|
\theta(\mathbf{x}-\mathbf{x}_0)
+
(1-\theta)(\mathbf{y}-\mathbf{x}_0)
\|
\leq
\theta
\|\mathbf{x}-\mathbf{x}_0\|
+
(1-\theta)
\|\mathbf{y}-\mathbf{x}_0\|.
\]

Thus norm constraints are fundamental convex constraints.

%%%%%%%%%%%%%%%%%%%%%%%%%%%%%%%%%%%%%%%%%%%%%%%%%%%%%%%%%%%%
\subsection{Convex Hull}
\label{subsec:convex-hull}
%%%%%%%%%%%%%%%%%%%%%%%%%%%%%%%%%%%%%%%%%%%%%%%%%%%%%%%%%%%%

\begin{definitionbox}{Convex Hull}
The convex hull of a set \(S\) is the smallest convex set containing
\(S\).
\end{definitionbox}

It consists of all finite convex combinations of points in \(S\):

\[
\boxed{
\operatorname{conv}(S)
=
\left\{
\sum_{i=1}^{k}
\theta_i\mathbf{x}_i:
\mathbf{x}_i\in S,
\;
\theta_i\geq0,
\;
\sum_i\theta_i=1
\right\}.
}
\]

%%%%%%%%%%%%%%%%%%%%%%%%%%%%%%%%%%%%%%%%%%%%%%%%%%%%%%%%%%%%
\subsection{Intersection of Convex Sets}
\label{subsec:intersection-convex-sets}
%%%%%%%%%%%%%%%%%%%%%%%%%%%%%%%%%%%%%%%%%%%%%%%%%%%%%%%%%%%%

The intersection of any collection of convex sets is convex.

If

\[
C_\alpha
\]

is convex for every index \(\alpha\), then

\[
\boxed{
\bigcap_\alpha C_\alpha
}
\]

is convex.

This is why systems of convex constraints preserve convexity of the
feasible region.

%%%%%%%%%%%%%%%%%%%%%%%%%%%%%%%%%%%%%%%%%%%%%%%%%%%%%%%%%%%%
\subsection{Convex Functions}
\label{subsec:convex-functions-convex-optimization}
%%%%%%%%%%%%%%%%%%%%%%%%%%%%%%%%%%%%%%%%%%%%%%%%%%%%%%%%%%%%

\begin{definitionbox}{Convex Function}
A function \(f:C\to\mathbb{R}\) on a convex domain \(C\) is convex if

\[
\boxed{
f
\left(
\theta\mathbf{x}
+
(1-\theta)\mathbf{y}
\right)
\leq
\theta f(\mathbf{x})
+
(1-\theta)f(\mathbf{y})
}
\]

for all

\[
\mathbf{x},\mathbf{y}\in C,
\qquad
0\leq\theta\leq1.
\]
\end{definitionbox}

%%%%%%%%%%%%%%%%%%%%%%%%%%%%%%%%%%%%%%%%%%%%%%%%%%%%%%%%%%%%
\subsection{Geometric Meaning of Convexity}
\label{subsec:geometric-meaning-convexity}
%%%%%%%%%%%%%%%%%%%%%%%%%%%%%%%%%%%%%%%%%%%%%%%%%%%%%%%%%%%%

For a convex function, the graph lies below the chord joining any two
points of the graph.

Thus,

\[
\boxed{
\text{chords lie above the function}.
}
\]

This geometric property is the source of many strong optimization
guarantees.

%%%%%%%%%%%%%%%%%%%%%%%%%%%%%%%%%%%%%%%%%%%%%%%%%%%%%%%%%%%%
\subsection{Strict Convexity}
\label{subsec:strict-convexity-convex-optimization}
%%%%%%%%%%%%%%%%%%%%%%%%%%%%%%%%%%%%%%%%%%%%%%%%%%%%%%%%%%%%

A function is strictly convex if

\[
\boxed{
f
\left(
\theta\mathbf{x}
+
(1-\theta)\mathbf{y}
\right)
<
\theta f(\mathbf{x})
+
(1-\theta)f(\mathbf{y})
}
\]

for distinct \(\mathbf{x}\neq\mathbf{y}\) and

\[
0<\theta<1.
\]

Strict convexity implies at most one minimizer over a convex set.

%%%%%%%%%%%%%%%%%%%%%%%%%%%%%%%%%%%%%%%%%%%%%%%%%%%%%%%%%%%%
\subsection{Strong Convexity}
\label{subsec:strong-convexity}
%%%%%%%%%%%%%%%%%%%%%%%%%%%%%%%%%%%%%%%%%%%%%%%%%%%%%%%%%%%%

A differentiable function is \(m\)-strongly convex if there exists

\[
m>0
\]

such that

\[
\boxed{
f(\mathbf{y})
\geq
f(\mathbf{x})
+
\nabla f(\mathbf{x})^T
(\mathbf{y}-\mathbf{x})
+
\frac{m}{2}
\|
\mathbf{y}-\mathbf{x}
\|_2^2.
}
\]

Strong convexity provides a uniform positive curvature condition.

%%%%%%%%%%%%%%%%%%%%%%%%%%%%%%%%%%%%%%%%%%%%%%%%%%%%%%%%%%%%
\subsection{Strong Convexity and Hessians}
\label{subsec:strong-convexity-hessian}
%%%%%%%%%%%%%%%%%%%%%%%%%%%%%%%%%%%%%%%%%%%%%%%%%%%%%%%%%%%%

If \(f\) is twice differentiable, a sufficient condition for
\(m\)-strong convexity is

\[
\boxed{
\nabla^2f(\mathbf{x})
\succeq
mI
}
\]

for every \(\mathbf{x}\) in the domain.

This means every Hessian eigenvalue is at least \(m\).

%%%%%%%%%%%%%%%%%%%%%%%%%%%%%%%%%%%%%%%%%%%%%%%%%%%%%%%%%%%%
\subsection{First-Order Convexity Criterion}
\label{subsec:first-order-convexity-criterion}
%%%%%%%%%%%%%%%%%%%%%%%%%%%%%%%%%%%%%%%%%%%%%%%%%%%%%%%%%%%%

For differentiable \(f\), convexity is equivalent to

\[
\boxed{
f(\mathbf{y})
\geq
f(\mathbf{x})
+
\nabla f(\mathbf{x})^T
(
\mathbf{y}-\mathbf{x}
)
}
\]

for all \(\mathbf{x},\mathbf{y}\) in the convex domain.

The tangent hyperplane globally underestimates a convex function.

%%%%%%%%%%%%%%%%%%%%%%%%%%%%%%%%%%%%%%%%%%%%%%%%%%%%%%%%%%%%
\subsection{Global Optimality from the Gradient}
\label{subsec:convex-gradient-global-optimality}
%%%%%%%%%%%%%%%%%%%%%%%%%%%%%%%%%%%%%%%%%%%%%%%%%%%%%%%%%%%%

Suppose \(f\) is differentiable and convex and

\[
\boxed{
\nabla f(\mathbf{x}^*)=0.
}
\]

Then, by the first-order convexity inequality,

\[
f(\mathbf{y})
\geq
f(\mathbf{x}^*)
\]

for every \(\mathbf{y}\).

Therefore,

\[
\boxed{
\mathbf{x}^*
\text{ is a global minimizer}.
}
\]

%%%%%%%%%%%%%%%%%%%%%%%%%%%%%%%%%%%%%%%%%%%%%%%%%%%%%%%%%%%%
\subsection{Worked Example: Global Convex Minimum}
\label{subsec:worked-global-convex-minimum}
%%%%%%%%%%%%%%%%%%%%%%%%%%%%%%%%%%%%%%%%%%%%%%%%%%%%%%%%%%%%

\begin{workedexamplebox}{Quadratic Convex Function}

Consider

\[
f(x,y)
=
x^2+2y^2-4x+8y.
\]

The gradient is

\[
\nabla f
=
\begin{pmatrix}
2x-4\\
4y+8
\end{pmatrix}.
\]

Set

\[
\nabla f=0.
\]

Then

\[
x=2,
\qquad
y=-2.
\]

The Hessian is

\[
\nabla^2f
=
\begin{pmatrix}
2&0\\
0&4
\end{pmatrix}.
\]

Since

\[
\nabla^2f\succ0,
\]

the function is strictly convex.

Therefore,

\begin{answerbox}
\[
\boxed{
(2,-2)
}
\]

is the unique global minimizer.
\end{answerbox}

\end{workedexamplebox}

%%%%%%%%%%%%%%%%%%%%%%%%%%%%%%%%%%%%%%%%%%%%%%%%%%%%%%%%%%%%
\subsection{Second-Order Convexity Criterion}
\label{subsec:second-order-convexity-criterion}
%%%%%%%%%%%%%%%%%%%%%%%%%%%%%%%%%%%%%%%%%%%%%%%%%%%%%%%%%%%%

If \(f\) is twice differentiable on a convex domain, then

\[
\boxed{
\nabla^2f(\mathbf{x})
\succeq0
}
\]

for every \(\mathbf{x}\) implies \(f\) is convex.

If

\[
\boxed{
\nabla^2f(\mathbf{x})
\succ0
}
\]

for every \(\mathbf{x}\), then \(f\) is strictly convex.

Positive definiteness everywhere is sufficient for strict convexity, but
strict convexity can hold even when the Hessian is singular at isolated
points.

%%%%%%%%%%%%%%%%%%%%%%%%%%%%%%%%%%%%%%%%%%%%%%%%%%%%%%%%%%%%
\subsection{Examples of Convex Functions}
\label{subsec:examples-convex-functions}
%%%%%%%%%%%%%%%%%%%%%%%%%%%%%%%%%%%%%%%%%%%%%%%%%%%%%%%%%%%%

Common convex functions include:

\[
\boxed{
f(x)=x^2,
}
\]

\[
\boxed{
f(x)=e^x,
}
\]

\[
\boxed{
f(x)=-\log x,
\qquad
x>0,
}
\]

\[
\boxed{
f(\mathbf{x})=\|\mathbf{x}\|,
}
\]

and

\[
\boxed{
f(\mathbf{x})
=
\frac12
\mathbf{x}^TQ\mathbf{x}
}
\]

when

\[
Q\succeq0.
\]

%%%%%%%%%%%%%%%%%%%%%%%%%%%%%%%%%%%%%%%%%%%%%%%%%%%%%%%%%%%%
\subsection{Epigraph}
\label{subsec:epigraph}
%%%%%%%%%%%%%%%%%%%%%%%%%%%%%%%%%%%%%%%%%%%%%%%%%%%%%%%%%%%%

\begin{definitionbox}{Epigraph}
The epigraph of \(f\) is

\[
\boxed{
\operatorname{epi}f
=
\{
(\mathbf{x},t):
f(\mathbf{x})\leq t
\}.
}
\]
\end{definitionbox}

A function is convex if and only if its epigraph is a convex set.

%%%%%%%%%%%%%%%%%%%%%%%%%%%%%%%%%%%%%%%%%%%%%%%%%%%%%%%%%%%%
\subsection{Epigraph Form of Optimization}
\label{subsec:epigraph-form-optimization}
%%%%%%%%%%%%%%%%%%%%%%%%%%%%%%%%%%%%%%%%%%%%%%%%%%%%%%%%%%%%

The problem

\[
\min_{\mathbf{x}}
f(\mathbf{x})
\]

can be rewritten as

\[
\boxed{
\begin{aligned}
\min_{\mathbf{x},t}
\quad&
t
\\
\text{subject to}
\quad&
f(\mathbf{x})\leq t.
\end{aligned}
}
\]

This is called epigraph form.

It converts a nonlinear objective into a linear objective with an
additional constraint.

%%%%%%%%%%%%%%%%%%%%%%%%%%%%%%%%%%%%%%%%%%%%%%%%%%%%%%%%%%%%
\subsection{Sublevel Sets}
\label{subsec:sublevel-sets}
%%%%%%%%%%%%%%%%%%%%%%%%%%%%%%%%%%%%%%%%%%%%%%%%%%%%%%%%%%%%

The \(\alpha\)-sublevel set of \(f\) is

\[
\boxed{
C_\alpha
=
\{
\mathbf{x}:
f(\mathbf{x})\leq\alpha
\}.
}
\]

If \(f\) is convex, every sublevel set is convex.

The converse is not generally true; functions whose sublevel sets are
convex are called quasiconvex.

%%%%%%%%%%%%%%%%%%%%%%%%%%%%%%%%%%%%%%%%%%%%%%%%%%%%%%%%%%%%
\subsection{Preserving Convexity by Nonnegative Sums}
\label{subsec:convex-nonnegative-sums}
%%%%%%%%%%%%%%%%%%%%%%%%%%%%%%%%%%%%%%%%%%%%%%%%%%%%%%%%%%%%

If \(f_1,\ldots,f_k\) are convex and

\[
\alpha_i\geq0,
\]

then

\[
\boxed{
f
=
\sum_{i=1}^{k}
\alpha_if_i
}
\]

is convex.

This allows complex convex objectives to be constructed from simpler
ones.

%%%%%%%%%%%%%%%%%%%%%%%%%%%%%%%%%%%%%%%%%%%%%%%%%%%%%%%%%%%%
\subsection{Affine Composition}
\label{subsec:convex-affine-composition}
%%%%%%%%%%%%%%%%%%%%%%%%%%%%%%%%%%%%%%%%%%%%%%%%%%%%%%%%%%%%

If \(f\) is convex, then

\[
\boxed{
g(\mathbf{x})
=
f(A\mathbf{x}+\mathbf{b})
}
\]

is convex on its domain.

Affine transformations therefore preserve convexity.

%%%%%%%%%%%%%%%%%%%%%%%%%%%%%%%%%%%%%%%%%%%%%%%%%%%%%%%%%%%%
\subsection{Pointwise Maximum}
\label{subsec:pointwise-maximum-convex}
%%%%%%%%%%%%%%%%%%%%%%%%%%%%%%%%%%%%%%%%%%%%%%%%%%%%%%%%%%%%

If

\[
f_1,\ldots,f_k
\]

are convex, then

\[
\boxed{
f(\mathbf{x})
=
\max_{i}
f_i(\mathbf{x})
}
\]

is convex.

More generally, the pointwise supremum of convex functions is convex when
finite on the domain of interest.

%%%%%%%%%%%%%%%%%%%%%%%%%%%%%%%%%%%%%%%%%%%%%%%%%%%%%%%%%%%%
\subsection{Worked Example: Maximum of Affine Functions}
\label{subsec:worked-max-affine-convex}
%%%%%%%%%%%%%%%%%%%%%%%%%%%%%%%%%%%%%%%%%%%%%%%%%%%%%%%%%%%%

\begin{workedexamplebox}{Piecewise Linear Convex Function}

Consider

\[
f(x)
=
\max
\{
2x+1,
-x+3
\}.
\]

Both component functions are affine and therefore convex.

The pointwise maximum of convex functions is convex.

Hence,

\begin{answerbox}
\[
\boxed{
f(x)
\text{ is convex}.
}
\]
\end{answerbox}

\end{workedexamplebox}

%%%%%%%%%%%%%%%%%%%%%%%%%%%%%%%%%%%%%%%%%%%%%%%%%%%%%%%%%%%%
\subsection{Composition Rules}
\label{subsec:convex-composition-rules}
%%%%%%%%%%%%%%%%%%%%%%%%%%%%%%%%%%%%%%%%%%%%%%%%%%%%%%%%%%%%

Suppose \(h\) is convex and nondecreasing in each argument, and
\(g_1,\ldots,g_k\) are convex.

Then under appropriate domain conditions,

\[
\boxed{
f(\mathbf{x})
=
h
(
g_1(\mathbf{x}),
\ldots,
g_k(\mathbf{x})
)
}
\]

is convex.

Composition rules are powerful but require monotonicity conditions to be
checked carefully.

%%%%%%%%%%%%%%%%%%%%%%%%%%%%%%%%%%%%%%%%%%%%%%%%%%%%%%%%%%%%
\subsection{Perspective Function}
\label{subsec:perspective-function}
%%%%%%%%%%%%%%%%%%%%%%%%%%%%%%%%%%%%%%%%%%%%%%%%%%%%%%%%%%%%

If

\[
f:\mathbb{R}^n\to\mathbb{R}
\]

is convex, its perspective is

\[
\boxed{
g(\mathbf{x},t)
=
t
f
\left(
\frac{\mathbf{x}}{t}
\right),
\qquad
t>0.
}
\]

The perspective of a convex function is convex.

This construction appears frequently in conic and statistical models.

%%%%%%%%%%%%%%%%%%%%%%%%%%%%%%%%%%%%%%%%%%%%%%%%%%%%%%%%%%%%
\subsection{Convex Cones}
\label{subsec:convex-cones}
%%%%%%%%%%%%%%%%%%%%%%%%%%%%%%%%%%%%%%%%%%%%%%%%%%%%%%%%%%%%

\begin{definitionbox}{Convex Cone}
A set \(K\) is a convex cone if for every
\(\mathbf{x},\mathbf{y}\in K\) and every
\(\alpha,\beta\geq0\),

\[
\boxed{
\alpha\mathbf{x}
+
\beta\mathbf{y}
\in K.
}
\]
\end{definitionbox}

Examples include:

\[
\boxed{
\mathbb{R}_+^n,
}
\]

and the cone of positive-semidefinite matrices.

%%%%%%%%%%%%%%%%%%%%%%%%%%%%%%%%%%%%%%%%%%%%%%%%%%%%%%%%%%%%
\subsection{Dual Cone}
\label{subsec:dual-cone}
%%%%%%%%%%%%%%%%%%%%%%%%%%%%%%%%%%%%%%%%%%%%%%%%%%%%%%%%%%%%

For a cone \(K\), the dual cone is

\[
\boxed{
K^*
=
\{
\mathbf{y}:
\mathbf{y}^T\mathbf{x}\geq0
\text{ for every }
\mathbf{x}\in K
\}.
}
\]

For the nonnegative orthant,

\[
\boxed{
(\mathbb{R}_+^n)^*
=
\mathbb{R}_+^n.
}
\]

Such a cone is self-dual.

%%%%%%%%%%%%%%%%%%%%%%%%%%%%%%%%%%%%%%%%%%%%%%%%%%%%%%%%%%%%
\subsection{Second-Order Cone}
\label{subsec:second-order-cone}
%%%%%%%%%%%%%%%%%%%%%%%%%%%%%%%%%%%%%%%%%%%%%%%%%%%%%%%%%%%%

The second-order cone, also called the Lorentz cone, is

\[
\boxed{
\mathcal{Q}^{n+1}
=
\left\{
(t,\mathbf{x}):
\|\mathbf{x}\|_2\leq t
\right\}.
}
\]

It is a convex cone.

Problems involving Euclidean norm constraints can often be written using
second-order cones.

%%%%%%%%%%%%%%%%%%%%%%%%%%%%%%%%%%%%%%%%%%%%%%%%%%%%%%%%%%%%
\subsection{Second-Order Cone Programming}
\label{subsec:second-order-cone-programming}
%%%%%%%%%%%%%%%%%%%%%%%%%%%%%%%%%%%%%%%%%%%%%%%%%%%%%%%%%%%%

A second-order cone program, abbreviated SOCP, has constraints of the form

\[
\boxed{
\|
A_i\mathbf{x}
+
\mathbf{b}_i
\|_2
\leq
\mathbf{c}_i^T\mathbf{x}
+
d_i.
}
\]

A typical SOCP has a linear objective and affine equality constraints.

SOCP includes linear programming as a special case.

%%%%%%%%%%%%%%%%%%%%%%%%%%%%%%%%%%%%%%%%%%%%%%%%%%%%%%%%%%%%
\subsection{Worked Example: Norm Constraint}
\label{subsec:worked-socp-norm}
%%%%%%%%%%%%%%%%%%%%%%%%%%%%%%%%%%%%%%%%%%%%%%%%%%%%%%%%%%%%

\begin{workedexamplebox}{Minimum Distance with a Linear Constraint}

Consider

\[
\min_{\mathbf{x}}
\|\mathbf{x}\|_2
\]

subject to

\[
A\mathbf{x}=\mathbf{b}.
\]

Introduce \(t\) and write

\[
\|\mathbf{x}\|_2\leq t.
\]

Then the problem becomes

\begin{answerbox}
\[
\boxed{
\begin{aligned}
\min_{\mathbf{x},t}
\quad&
t
\\
\text{subject to}
\quad&
\|\mathbf{x}\|_2\leq t,
\\
&
A\mathbf{x}=\mathbf{b}.
\end{aligned}
}
\]
\end{answerbox}

This is an SOCP formulation.

\end{workedexamplebox}

%%%%%%%%%%%%%%%%%%%%%%%%%%%%%%%%%%%%%%%%%%%%%%%%%%%%%%%%%%%%
\subsection{Positive-Semidefinite Cone}
\label{subsec:positive-semidefinite-cone}
%%%%%%%%%%%%%%%%%%%%%%%%%%%%%%%%%%%%%%%%%%%%%%%%%%%%%%%%%%%%

Let \(\mathbb{S}^n\) denote the symmetric \(n\times n\) matrices.

The positive-semidefinite cone is

\[
\boxed{
\mathbb{S}_+^n
=
\{
X\in\mathbb{S}^n:
X\succeq0
\}.
}
\]

It is convex.

%%%%%%%%%%%%%%%%%%%%%%%%%%%%%%%%%%%%%%%%%%%%%%%%%%%%%%%%%%%%
\subsection{Semidefinite Programming}
\label{subsec:semidefinite-programming}
%%%%%%%%%%%%%%%%%%%%%%%%%%%%%%%%%%%%%%%%%%%%%%%%%%%%%%%%%%%%

A semidefinite program, abbreviated SDP, has a linear objective and
positive-semidefinite matrix constraints.

A standard form is

\[
\boxed{
\begin{aligned}
\min_X
\quad&
\langle C,X\rangle
\\
\text{subject to}
\quad&
\langle A_i,X\rangle=b_i,
\qquad
i=1,\ldots,m,
\\
&
X\succeq0.
\end{aligned}
}
\]

Here

\[
\boxed{
\langle A,X\rangle
=
\operatorname{tr}(A^TX).
}
\]

%%%%%%%%%%%%%%%%%%%%%%%%%%%%%%%%%%%%%%%%%%%%%%%%%%%%%%%%%%%%
\subsection{Why Semidefinite Programming Matters}
\label{subsec:why-sdp-matters}
%%%%%%%%%%%%%%%%%%%%%%%%%%%%%%%%%%%%%%%%%%%%%%%%%%%%%%%%%%%%

SDPs can represent or relax problems involving:

\[
\boxed{
\text{matrix inequalities},
}
\]

\[
\boxed{
\text{control stability},
}
\]

\[
\boxed{
\text{quadratic optimization},
}
\]

\[
\boxed{
\text{combinatorial optimization},
}
\]

and

\[
\boxed{
\text{moment problems}.
}
\]

They form one of the major classes of conic optimization.

%%%%%%%%%%%%%%%%%%%%%%%%%%%%%%%%%%%%%%%%%%%%%%%%%%%%%%%%%%%%
\subsection{Convex Quadratic Programming}
\label{subsec:convex-quadratic-programming}
%%%%%%%%%%%%%%%%%%%%%%%%%%%%%%%%%%%%%%%%%%%%%%%%%%%%%%%%%%%%

A convex quadratic program has form

\[
\boxed{
\begin{aligned}
\min_{\mathbf{x}}
\quad&
\frac12
\mathbf{x}^TQ\mathbf{x}
+
\mathbf{c}^T\mathbf{x}
\\
\text{subject to}
\quad&
A\mathbf{x}\leq\mathbf{b},
\\
&
E\mathbf{x}=\mathbf{d},
\end{aligned}
}
\]

where

\[
\boxed{
Q\succeq0.
}
\]

If

\[
Q\succ0,
\]

the objective is strictly convex.

%%%%%%%%%%%%%%%%%%%%%%%%%%%%%%%%%%%%%%%%%%%%%%%%%%%%%%%%%%%%
\subsection{Worked Example: Convex Quadratic Program}
\label{subsec:worked-convex-qp}
%%%%%%%%%%%%%%%%%%%%%%%%%%%%%%%%%%%%%%%%%%%%%%%%%%%%%%%%%%%%

\begin{workedexamplebox}{Projection onto a Half-Space}

Minimize

\[
\frac12
\left[
(x-3)^2
+
(y-2)^2
\right]
\]

subject to

\[
x+y\leq2.
\]

This finds the Euclidean projection of \((3,2)\) onto the half-space.

Write

\[
g(x,y)=x+y-2\leq0.
\]

The Lagrangian is

\[
\mathcal{L}
=
\frac12
\left[
(x-3)^2+(y-2)^2
\right]
+
\mu(x+y-2).
\]

Stationarity gives

\[
x-3+\mu=0,
\]

\[
y-2+\mu=0.
\]

Thus,

\[
x=3-\mu,
\]

\[
y=2-\mu.
\]

The unconstrained point \((3,2)\) is infeasible, so the constraint is
active:

\[
x+y=2.
\]

Therefore,

\[
5-2\mu=2,
\]

so

\[
\mu=\frac32.
\]

Hence,

\[
x=\frac32,
\qquad
y=\frac12.
\]

Thus,

\begin{answerbox}
\[
\boxed{
(x^*,y^*)
=
\left(
\frac32,
\frac12
\right).
}
\]
\end{answerbox}

\end{workedexamplebox}

%%%%%%%%%%%%%%%%%%%%%%%%%%%%%%%%%%%%%%%%%%%%%%%%%%%%%%%%%%%%
\subsection{Projection onto a Convex Set}
\label{subsec:projection-convex-set}
%%%%%%%%%%%%%%%%%%%%%%%%%%%%%%%%%%%%%%%%%%%%%%%%%%%%%%%%%%%%

For a closed convex set \(C\), the Euclidean projection of
\(\mathbf{x}\) onto \(C\) is

\[
\boxed{
P_C(\mathbf{x})
=
\operatorname*{arg\,min}_{\mathbf{y}\in C}
\frac12
\|
\mathbf{x}-\mathbf{y}
\|_2^2.
}
\]

For nonempty closed convex \(C\), this projection exists and is unique.

%%%%%%%%%%%%%%%%%%%%%%%%%%%%%%%%%%%%%%%%%%%%%%%%%%%%%%%%%%%%
\subsection{Projection Characterization}
\label{subsec:projection-characterization}
%%%%%%%%%%%%%%%%%%%%%%%%%%%%%%%%%%%%%%%%%%%%%%%%%%%%%%%%%%%%

A point

\[
\mathbf{y}=P_C(\mathbf{x})
\]

if and only if

\[
\boxed{
(\mathbf{x}-\mathbf{y})^T
(\mathbf{z}-\mathbf{y})
\leq0
}
\]

for every

\[
\mathbf{z}\in C.
\]

Geometrically, the vector from the projection to the original point forms
an obtuse or right angle with every feasible direction.

%%%%%%%%%%%%%%%%%%%%%%%%%%%%%%%%%%%%%%%%%%%%%%%%%%%%%%%%%%%%
\subsection{Supporting Hyperplanes}
\label{subsec:supporting-hyperplanes}
%%%%%%%%%%%%%%%%%%%%%%%%%%%%%%%%%%%%%%%%%%%%%%%%%%%%%%%%%%%%

A hyperplane

\[
\boxed{
\mathbf{a}^T\mathbf{x}=b
}
\]

supports a convex set \(C\) at \(\mathbf{x}_0\in C\) if

\[
\boxed{
\mathbf{a}^T\mathbf{x}
\leq b
}
\]

for every \(\mathbf{x}\in C\), while

\[
\boxed{
\mathbf{a}^T\mathbf{x}_0=b.
}
\]

Supporting hyperplanes generalize tangent lines to convex sets.

%%%%%%%%%%%%%%%%%%%%%%%%%%%%%%%%%%%%%%%%%%%%%%%%%%%%%%%%%%%%
\subsection{Separating Hyperplanes}
\label{subsec:separating-hyperplanes}
%%%%%%%%%%%%%%%%%%%%%%%%%%%%%%%%%%%%%%%%%%%%%%%%%%%%%%%%%%%%

If two convex sets do not intersect under suitable regularity
conditions, a hyperplane can often separate them.

This geometric principle underlies:

\[
\boxed{
\text{duality},
}
\]

\[
\boxed{
\text{optimality certificates},
}
\]

and

\[
\boxed{
\text{alternative theorems}.
}
\]

%%%%%%%%%%%%%%%%%%%%%%%%%%%%%%%%%%%%%%%%%%%%%%%%%%%%%%%%%%%%
\subsection{Subgradients}
\label{subsec:subgradients}
%%%%%%%%%%%%%%%%%%%%%%%%%%%%%%%%%%%%%%%%%%%%%%%%%%%%%%%%%%%%

For a convex function \(f\), a vector \(\mathbf{g}\) is a subgradient at
\(\mathbf{x}\) if

\[
\boxed{
f(\mathbf{y})
\geq
f(\mathbf{x})
+
\mathbf{g}^T
(
\mathbf{y}-\mathbf{x}
)
}
\]

for every \(\mathbf{y}\).

The set of all subgradients is the subdifferential:

\[
\boxed{
\partial f(\mathbf{x}).
}
\]

%%%%%%%%%%%%%%%%%%%%%%%%%%%%%%%%%%%%%%%%%%%%%%%%%%%%%%%%%%%%
\subsection{Differentiable Case}
\label{subsec:subgradient-differentiable-case}
%%%%%%%%%%%%%%%%%%%%%%%%%%%%%%%%%%%%%%%%%%%%%%%%%%%%%%%%%%%%

If \(f\) is convex and differentiable at \(\mathbf{x}\), then

\[
\boxed{
\partial f(\mathbf{x})
=
\{
\nabla f(\mathbf{x})
\}.
}
\]

Thus subgradients extend gradients to nonsmooth convex functions.

%%%%%%%%%%%%%%%%%%%%%%%%%%%%%%%%%%%%%%%%%%%%%%%%%%%%%%%%%%%%
\subsection{Subgradient of Absolute Value}
\label{subsec:subgradient-absolute-value}
%%%%%%%%%%%%%%%%%%%%%%%%%%%%%%%%%%%%%%%%%%%%%%%%%%%%%%%%%%%%

For

\[
f(x)=|x|,
\]

we have

\[
\boxed{
\partial f(x)
=
\begin{cases}
\{1\},
&
x>0,
\\
[-1,1],
&
x=0,
\\
\{-1\},
&
x<0.
\end{cases}
}
\]

The interval at \(x=0\) replaces the missing ordinary derivative.

%%%%%%%%%%%%%%%%%%%%%%%%%%%%%%%%%%%%%%%%%%%%%%%%%%%%%%%%%%%%
\subsection{Subgradient Optimality}
\label{subsec:subgradient-optimality}
%%%%%%%%%%%%%%%%%%%%%%%%%%%%%%%%%%%%%%%%%%%%%%%%%%%%%%%%%%%%

For a convex function \(f\),

\[
\boxed{
\mathbf{x}^*
\text{ minimizes }f
\iff
\mathbf{0}
\in
\partial f(\mathbf{x}^*).
}
\]

This is the nonsmooth analogue of

\[
\nabla f(\mathbf{x}^*)=0.
\]

%%%%%%%%%%%%%%%%%%%%%%%%%%%%%%%%%%%%%%%%%%%%%%%%%%%%%%%%%%%%
\subsection{Worked Example: Nonsmooth Convex Minimum}
\label{subsec:worked-subgradient-minimum}
%%%%%%%%%%%%%%%%%%%%%%%%%%%%%%%%%%%%%%%%%%%%%%%%%%%%%%%%%%%%

\begin{workedexamplebox}{Absolute Value}

Consider

\[
f(x)=|x-2|.
\]

At

\[
x=2,
\]

the subdifferential is

\[
\partial f(2)
=
[-1,1].
\]

Since

\[
0\in[-1,1],
\]

we conclude

\begin{answerbox}
\[
\boxed{
x^*=2
}
\]

is the global minimizer.
\end{answerbox}

\end{workedexamplebox}

%%%%%%%%%%%%%%%%%%%%%%%%%%%%%%%%%%%%%%%%%%%%%%%%%%%%%%%%%%%%
\subsection{Normal Cone}
\label{subsec:normal-cone-convex}
%%%%%%%%%%%%%%%%%%%%%%%%%%%%%%%%%%%%%%%%%%%%%%%%%%%%%%%%%%%%

For a convex set \(C\) and \(\mathbf{x}\in C\), the normal cone is

\[
\boxed{
N_C(\mathbf{x})
=
\{
\mathbf{v}:
\mathbf{v}^T
(
\mathbf{y}-\mathbf{x}
)
\leq0
\text{ for all }
\mathbf{y}\in C
\}.
}
\]

At an interior point,

\[
\boxed{
N_C(\mathbf{x})=\{\mathbf{0}\}.
}
\]

At a boundary point, it contains outward normal directions.

%%%%%%%%%%%%%%%%%%%%%%%%%%%%%%%%%%%%%%%%%%%%%%%%%%%%%%%%%%%%
\subsection{Constrained Convex Optimality}
\label{subsec:constrained-convex-optimality}
%%%%%%%%%%%%%%%%%%%%%%%%%%%%%%%%%%%%%%%%%%%%%%%%%%%%%%%%%%%%

For differentiable convex \(f\) over a closed convex set \(C\),
\(\mathbf{x}^*\) is optimal if and only if

\[
\boxed{
-\nabla f(\mathbf{x}^*)
\in
N_C(\mathbf{x}^*).
}
\]

Equivalently,

\[
\boxed{
\nabla f(\mathbf{x}^*)^T
(
\mathbf{x}-\mathbf{x}^*
)
\geq0
}
\]

for every

\[
\mathbf{x}\in C.
\]

%%%%%%%%%%%%%%%%%%%%%%%%%%%%%%%%%%%%%%%%%%%%%%%%%%%%%%%%%%%%
\subsection{Lagrange Dual Function}
\label{subsec:convex-lagrange-dual-function}
%%%%%%%%%%%%%%%%%%%%%%%%%%%%%%%%%%%%%%%%%%%%%%%%%%%%%%%%%%%%

Consider the primal problem

\[
\boxed{
\begin{aligned}
\min_{\mathbf{x}}
\quad&
f_0(\mathbf{x})
\\
\text{subject to}
\quad&
f_i(\mathbf{x})\leq0,
\\
&
h_j(\mathbf{x})=0.
\end{aligned}
}
\]

The Lagrangian is

\[
\boxed{
\mathcal{L}
(
\mathbf{x},
\boldsymbol{\lambda},
\boldsymbol{\nu}
)
=
f_0(\mathbf{x})
+
\sum_{i=1}^{m}
\lambda_if_i(\mathbf{x})
+
\sum_{j=1}^{p}
\nu_jh_j(\mathbf{x}),
}
\]

with

\[
\lambda_i\geq0.
\]

The dual function is

\[
\boxed{
q
(
\boldsymbol{\lambda},
\boldsymbol{\nu}
)
=
\inf_{\mathbf{x}}
\mathcal{L}
(
\mathbf{x},
\boldsymbol{\lambda},
\boldsymbol{\nu}
).
}
\]

%%%%%%%%%%%%%%%%%%%%%%%%%%%%%%%%%%%%%%%%%%%%%%%%%%%%%%%%%%%%
\subsection{Dual Function Gives Lower Bounds}
\label{subsec:dual-function-lower-bound}
%%%%%%%%%%%%%%%%%%%%%%%%%%%%%%%%%%%%%%%%%%%%%%%%%%%%%%%%%%%%

Suppose \(\mathbf{x}\) is primal feasible and
\(\boldsymbol{\lambda}\geq0\).

Then

\[
f_i(\mathbf{x})\leq0,
\]

so

\[
\sum_i
\lambda_if_i(\mathbf{x})
\leq0.
\]

Also,

\[
h_j(\mathbf{x})=0.
\]

Therefore,

\[
\mathcal{L}
(
\mathbf{x},
\boldsymbol{\lambda},
\boldsymbol{\nu}
)
\leq
f_0(\mathbf{x}).
\]

Since the dual function is the infimum over \(\mathbf{x}\),

\[
\boxed{
q
(
\boldsymbol{\lambda},
\boldsymbol{\nu}
)
\leq
f_0(\mathbf{x})
}
\]

for every primal feasible \(\mathbf{x}\).

%%%%%%%%%%%%%%%%%%%%%%%%%%%%%%%%%%%%%%%%%%%%%%%%%%%%%%%%%%%%
\subsection{Lagrange Dual Problem}
\label{subsec:lagrange-dual-problem}
%%%%%%%%%%%%%%%%%%%%%%%%%%%%%%%%%%%%%%%%%%%%%%%%%%%%%%%%%%%%

Because the dual function produces lower bounds on the primal minimum,
the dual problem maximizes that bound:

\[
\boxed{
\begin{aligned}
\max_{\boldsymbol{\lambda},\boldsymbol{\nu}}
\quad&
q
(
\boldsymbol{\lambda},
\boldsymbol{\nu}
)
\\
\text{subject to}
\quad&
\boldsymbol{\lambda}\geq0.
\end{aligned}
}
\]

The dual problem is always a convex optimization problem because the dual
function is concave in its dual variables.

%%%%%%%%%%%%%%%%%%%%%%%%%%%%%%%%%%%%%%%%%%%%%%%%%%%%%%%%%%%%
\subsection{Weak Duality}
\label{subsec:convex-weak-duality}
%%%%%%%%%%%%%%%%%%%%%%%%%%%%%%%%%%%%%%%%%%%%%%%%%%%%%%%%%%%%

Let

\[
p^*
\]

be the primal optimal value and

\[
d^*
\]

the dual optimal value.

Then

\[
\boxed{
d^*\leq p^*.
}
\]

This is weak duality.

The difference

\[
\boxed{
p^*-d^*
}
\]

is the duality gap.

%%%%%%%%%%%%%%%%%%%%%%%%%%%%%%%%%%%%%%%%%%%%%%%%%%%%%%%%%%%%
\subsection{Strong Duality}
\label{subsec:convex-strong-duality}
%%%%%%%%%%%%%%%%%%%%%%%%%%%%%%%%%%%%%%%%%%%%%%%%%%%%%%%%%%%%

Strong duality means

\[
\boxed{
d^*=p^*.
}
\]

For convex problems, strong duality holds under important regularity
conditions.

One of the most useful is Slater's condition.

%%%%%%%%%%%%%%%%%%%%%%%%%%%%%%%%%%%%%%%%%%%%%%%%%%%%%%%%%%%%
\subsection{Slater's Condition}
\label{subsec:slaters-condition}
%%%%%%%%%%%%%%%%%%%%%%%%%%%%%%%%%%%%%%%%%%%%%%%%%%%%%%%%%%%%

For a convex problem with affine equalities, Slater's condition requires
existence of a point satisfying

\[
\boxed{
f_i(\mathbf{x})<0
}
\]

for every nonlinear inequality constraint, together with all affine
equalities.

Such a point is called strictly feasible.

Under standard assumptions, Slater's condition implies strong duality.

%%%%%%%%%%%%%%%%%%%%%%%%%%%%%%%%%%%%%%%%%%%%%%%%%%%%%%%%%%%%
\subsection{KKT Sufficiency in Convex Optimization}
\label{subsec:kkt-sufficiency-convex}
%%%%%%%%%%%%%%%%%%%%%%%%%%%%%%%%%%%%%%%%%%%%%%%%%%%%%%%%%%%%

For a convex optimization problem, suppose:

\[
f_0
\]

and all inequality functions \(f_i\) are convex,

\[
h_j
\]

are affine, and KKT conditions hold.

Then the KKT point is globally optimal.

Under Slater's condition, the KKT conditions are also necessary for
optimality.

%%%%%%%%%%%%%%%%%%%%%%%%%%%%%%%%%%%%%%%%%%%%%%%%%%%%%%%%%%%%
\subsection{Convex KKT Conditions}
\label{subsec:convex-kkt-conditions}
%%%%%%%%%%%%%%%%%%%%%%%%%%%%%%%%%%%%%%%%%%%%%%%%%%%%%%%%%%%%

The conditions are:

\[
\boxed{
\nabla f_0(\mathbf{x}^*)
+
\sum_i
\lambda_i^*
\nabla f_i(\mathbf{x}^*)
+
\sum_j
\nu_j^*
\nabla h_j(\mathbf{x}^*)
=
0,
}
\]

\[
\boxed{
f_i(\mathbf{x}^*)\leq0,
}
\]

\[
\boxed{
h_j(\mathbf{x}^*)=0,
}
\]

\[
\boxed{
\lambda_i^*\geq0,
}
\]

and

\[
\boxed{
\lambda_i^*
f_i(\mathbf{x}^*)
=
0.
}
\]

%%%%%%%%%%%%%%%%%%%%%%%%%%%%%%%%%%%%%%%%%%%%%%%%%%%%%%%%%%%%
\subsection{Worked Example: Convex KKT Certificate}
\label{subsec:worked-convex-kkt}
%%%%%%%%%%%%%%%%%%%%%%%%%%%%%%%%%%%%%%%%%%%%%%%%%%%%%%%%%%%%

\begin{workedexamplebox}{Global Optimality from KKT}

Minimize

\[
f(x)=x^2
\]

subject to

\[
1-x\leq0.
\]

The problem is convex.

The Lagrangian is

\[
\mathcal{L}
=
x^2+\lambda(1-x).
\]

Stationarity gives

\[
2x-\lambda=0.
\]

The unconstrained minimizer \(x=0\) is infeasible, so the constraint is
active:

\[
x=1.
\]

Thus,

\[
\lambda=2.
\]

All KKT conditions hold.

Because the problem is convex,

\begin{answerbox}
\[
\boxed{
x^*=1
}
\]

is the global minimizer.
\end{answerbox}

\end{workedexamplebox}

%%%%%%%%%%%%%%%%%%%%%%%%%%%%%%%%%%%%%%%%%%%%%%%%%%%%%%%%%%%%
\subsection{Dual Interpretation of Multipliers}
\label{subsec:convex-dual-multiplier-interpretation}
%%%%%%%%%%%%%%%%%%%%%%%%%%%%%%%%%%%%%%%%%%%%%%%%%%%%%%%%%%%%

Dual variables quantify the value of relaxing constraints.

Suppose the \(i\)-th constraint becomes

\[
f_i(\mathbf{x})\leq u_i.
\]

Under suitable differentiability and regularity, the optimal multiplier
can describe local sensitivity of the optimal value with respect to
\(u_i\).

This generalizes shadow prices from linear programming.

%%%%%%%%%%%%%%%%%%%%%%%%%%%%%%%%%%%%%%%%%%%%%%%%%%%%%%%%%%%%
\subsection{Convex Conjugate}
\label{subsec:convex-conjugate}
%%%%%%%%%%%%%%%%%%%%%%%%%%%%%%%%%%%%%%%%%%%%%%%%%%%%%%%%%%%%

The convex conjugate of \(f\) is

\[
\boxed{
f^*(\mathbf{y})
=
\sup_{\mathbf{x}}
\left[
\mathbf{y}^T\mathbf{x}
-
f(\mathbf{x})
\right].
}
\]

This is also called the Fenchel conjugate.

The conjugate is always convex, even if the original function is not.

%%%%%%%%%%%%%%%%%%%%%%%%%%%%%%%%%%%%%%%%%%%%%%%%%%%%%%%%%%%%
\subsection{Fenchel--Young Inequality}
\label{subsec:fenchel-young}
%%%%%%%%%%%%%%%%%%%%%%%%%%%%%%%%%%%%%%%%%%%%%%%%%%%%%%%%%%%%

By the definition of the conjugate,

\[
f^*(\mathbf{y})
\geq
\mathbf{y}^T\mathbf{x}
-
f(\mathbf{x}).
\]

Therefore,

\[
\boxed{
f(\mathbf{x})
+
f^*(\mathbf{y})
\geq
\mathbf{x}^T\mathbf{y}.
}
\]

This is the Fenchel--Young inequality.

Equality holds precisely under an appropriate subgradient relation.

%%%%%%%%%%%%%%%%%%%%%%%%%%%%%%%%%%%%%%%%%%%%%%%%%%%%%%%%%%%%
\subsection{Fenchel Duality Preview}
\label{subsec:fenchel-duality-preview}
%%%%%%%%%%%%%%%%%%%%%%%%%%%%%%%%%%%%%%%%%%%%%%%%%%%%%%%%%%%%

Problems of the form

\[
\boxed{
\min_{\mathbf{x}}
f(\mathbf{x})
+
g(A\mathbf{x})
}
\]

can often be dualized using convex conjugates.

Fenchel duality provides a powerful framework extending ordinary
Lagrange duality.

It is particularly important in signal processing and machine learning.

%%%%%%%%%%%%%%%%%%%%%%%%%%%%%%%%%%%%%%%%%%%%%%%%%%%%%%%%%%%%
\subsection{Proximal Operator}
\label{subsec:proximal-operator}
%%%%%%%%%%%%%%%%%%%%%%%%%%%%%%%%%%%%%%%%%%%%%%%%%%%%%%%%%%%%

For a proper convex function \(f\), the proximal operator is

\[
\boxed{
\operatorname{prox}_{\lambda f}(\mathbf{v})
=
\operatorname*{arg\,min}_{\mathbf{x}}
\left[
f(\mathbf{x})
+
\frac{1}{2\lambda}
\|
\mathbf{x}-\mathbf{v}
\|_2^2
\right],
}
\]

where

\[
\lambda>0.
\]

The quadratic term keeps the solution near \(\mathbf{v}\).

%%%%%%%%%%%%%%%%%%%%%%%%%%%%%%%%%%%%%%%%%%%%%%%%%%%%%%%%%%%%
\subsection{Projection as a Proximal Operator}
\label{subsec:projection-proximal}
%%%%%%%%%%%%%%%%%%%%%%%%%%%%%%%%%%%%%%%%%%%%%%%%%%%%%%%%%%%%

Let the indicator function of a convex set \(C\) be

\[
\boxed{
I_C(\mathbf{x})
=
\begin{cases}
0,
&
\mathbf{x}\in C,
\\
+\infty,
&
\mathbf{x}\notin C.
\end{cases}
}
\]

Then

\[
\boxed{
\operatorname{prox}_{I_C}(\mathbf{v})
=
P_C(\mathbf{v}).
}
\]

Thus projection is a special proximal operation.

%%%%%%%%%%%%%%%%%%%%%%%%%%%%%%%%%%%%%%%%%%%%%%%%%%%%%%%%%%%%
\subsection{Soft Thresholding}
\label{subsec:soft-thresholding}
%%%%%%%%%%%%%%%%%%%%%%%%%%%%%%%%%%%%%%%%%%%%%%%%%%%%%%%%%%%%

For

\[
f(x)=|x|,
\]

the proximal operator is

\[
\boxed{
\operatorname{prox}_{\lambda|\cdot|}(v)
=
\operatorname{sign}(v)
\max
\{
|v|-\lambda,
0
\}.
}
\]

This is called soft thresholding.

It is fundamental in sparse optimization.

%%%%%%%%%%%%%%%%%%%%%%%%%%%%%%%%%%%%%%%%%%%%%%%%%%%%%%%%%%%%
\subsection{Gradient Descent for Convex Functions}
\label{subsec:gradient-descent-convex}
%%%%%%%%%%%%%%%%%%%%%%%%%%%%%%%%%%%%%%%%%%%%%%%%%%%%%%%%%%%%

For differentiable convex \(f\),

\[
\boxed{
\mathbf{x}_{k+1}
=
\mathbf{x}_k
-
\alpha_k
\nabla f(\mathbf{x}_k).
}
\]

If the gradient is Lipschitz continuous and the step size is chosen
appropriately, the method converges to a global minimizer under standard
assumptions.

%%%%%%%%%%%%%%%%%%%%%%%%%%%%%%%%%%%%%%%%%%%%%%%%%%%%%%%%%%%%
\subsection{Lipschitz Gradient}
\label{subsec:lipschitz-gradient}
%%%%%%%%%%%%%%%%%%%%%%%%%%%%%%%%%%%%%%%%%%%%%%%%%%%%%%%%%%%%

A gradient is \(L\)-Lipschitz if

\[
\boxed{
\|
\nabla f(\mathbf{x})
-
\nabla f(\mathbf{y})
\|_2
\leq
L
\|
\mathbf{x}-\mathbf{y}
\|_2.
}
\]

For twice-differentiable functions, a sufficient condition is

\[
\boxed{
\nabla^2f(\mathbf{x})
\preceq
LI.
}
\]

%%%%%%%%%%%%%%%%%%%%%%%%%%%%%%%%%%%%%%%%%%%%%%%%%%%%%%%%%%%%
\subsection{Descent Lemma}
\label{subsec:descent-lemma}
%%%%%%%%%%%%%%%%%%%%%%%%%%%%%%%%%%%%%%%%%%%%%%%%%%%%%%%%%%%%

If \(\nabla f\) is \(L\)-Lipschitz, then

\[
\boxed{
f(\mathbf{y})
\leq
f(\mathbf{x})
+
\nabla f(\mathbf{x})^T
(
\mathbf{y}-\mathbf{x}
)
+
\frac{L}{2}
\|
\mathbf{y}-\mathbf{x}
\|_2^2.
}
\]

This inequality is central to convergence analysis of gradient methods.

%%%%%%%%%%%%%%%%%%%%%%%%%%%%%%%%%%%%%%%%%%%%%%%%%%%%%%%%%%%%
\subsection{Fixed-Step Gradient Descent}
\label{subsec:fixed-step-gradient-descent}
%%%%%%%%%%%%%%%%%%%%%%%%%%%%%%%%%%%%%%%%%%%%%%%%%%%%%%%%%%%%

For an \(L\)-smooth convex function, a common choice is

\[
\boxed{
\alpha
=
\frac1L.
}
\]

Then

\[
\boxed{
\mathbf{x}_{k+1}
=
\mathbf{x}_k
-
\frac1L
\nabla f(\mathbf{x}_k).
}
\]

The objective converges sublinearly under standard assumptions.

%%%%%%%%%%%%%%%%%%%%%%%%%%%%%%%%%%%%%%%%%%%%%%%%%%%%%%%%%%%%
\subsection{Strongly Convex Convergence}
\label{subsec:strong-convex-gradient-convergence}
%%%%%%%%%%%%%%%%%%%%%%%%%%%%%%%%%%%%%%%%%%%%%%%%%%%%%%%%%%%%

If \(f\) is both:

\[
\boxed{
m\text{-strongly convex},
}
\]

and

\[
\boxed{
L\text{-smooth},
}
\]

gradient descent can converge linearly.

The ratio

\[
\boxed{
\kappa
=
\frac{L}{m}
}
\]

is a condition-number-like measure.

Large \(\kappa\) generally means slower convergence.

%%%%%%%%%%%%%%%%%%%%%%%%%%%%%%%%%%%%%%%%%%%%%%%%%%%%%%%%%%%%
\subsection{Subgradient Method}
\label{subsec:subgradient-method}
%%%%%%%%%%%%%%%%%%%%%%%%%%%%%%%%%%%%%%%%%%%%%%%%%%%%%%%%%%%%

For nonsmooth convex \(f\), choose

\[
\mathbf{g}_k
\in
\partial f(\mathbf{x}_k)
\]

and update

\[
\boxed{
\mathbf{x}_{k+1}
=
\mathbf{x}_k
-
\alpha_k
\mathbf{g}_k.
}
\]

The subgradient method is simple but typically converges more slowly than
gradient descent on smooth functions.

%%%%%%%%%%%%%%%%%%%%%%%%%%%%%%%%%%%%%%%%%%%%%%%%%%%%%%%%%%%%
\subsection{Projected Gradient Descent}
\label{subsec:projected-gradient-descent}
%%%%%%%%%%%%%%%%%%%%%%%%%%%%%%%%%%%%%%%%%%%%%%%%%%%%%%%%%%%%

For a closed convex feasible set \(C\),

\[
\boxed{
\mathbf{x}_{k+1}
=
P_C
\left(
\mathbf{x}_k
-
\alpha_k
\nabla f(\mathbf{x}_k)
\right).
}
\]

The gradient step reduces the objective locally, while projection restores
feasibility.

%%%%%%%%%%%%%%%%%%%%%%%%%%%%%%%%%%%%%%%%%%%%%%%%%%%%%%%%%%%%
\subsection{Proximal Gradient Method}
\label{subsec:proximal-gradient}
%%%%%%%%%%%%%%%%%%%%%%%%%%%%%%%%%%%%%%%%%%%%%%%%%%%%%%%%%%%%

Suppose

\[
\boxed{
F(\mathbf{x})
=
f(\mathbf{x})
+
g(\mathbf{x}),
}
\]

where \(f\) is smooth convex and \(g\) is convex but possibly nonsmooth.

The proximal-gradient update is

\[
\boxed{
\mathbf{x}_{k+1}
=
\operatorname{prox}_{\alpha g}
\left[
\mathbf{x}_k
-
\alpha
\nabla f(\mathbf{x}_k)
\right].
}
\]

This method is central to sparse and composite optimization.

%%%%%%%%%%%%%%%%%%%%%%%%%%%%%%%%%%%%%%%%%%%%%%%%%%%%%%%%%%%%
\subsection{Accelerated Gradient Methods Preview}
\label{subsec:accelerated-gradient-preview}
%%%%%%%%%%%%%%%%%%%%%%%%%%%%%%%%%%%%%%%%%%%%%%%%%%%%%%%%%%%%

Accelerated first-order methods combine gradient steps with momentum-like
extrapolation.

For smooth convex optimization, acceleration can improve the objective
convergence rate from roughly

\[
\boxed{
O\left(\frac1k\right)
}
\]

to

\[
\boxed{
O\left(\frac1{k^2}\right)
}
\]

under standard assumptions.

%%%%%%%%%%%%%%%%%%%%%%%%%%%%%%%%%%%%%%%%%%%%%%%%%%%%%%%%%%%%
\subsection{Coordinate Descent Preview}
\label{subsec:coordinate-descent-preview}
%%%%%%%%%%%%%%%%%%%%%%%%%%%%%%%%%%%%%%%%%%%%%%%%%%%%%%%%%%%%

Coordinate descent updates only one or a small group of variables at a
time.

A generic step changes

\[
\boxed{
x_i
}
\]

while holding other components fixed.

This can be effective for very large problems with simple coordinatewise
structure.

%%%%%%%%%%%%%%%%%%%%%%%%%%%%%%%%%%%%%%%%%%%%%%%%%%%%%%%%%%%%
\subsection{Interior-Point Methods}
\label{subsec:convex-interior-point}
%%%%%%%%%%%%%%%%%%%%%%%%%%%%%%%%%%%%%%%%%%%%%%%%%%%%%%%%%%%%

For a convex problem with inequalities

\[
f_i(\mathbf{x})\leq0,
\]

replace the constraints temporarily with the logarithmic barrier

\[
\boxed{
\phi(\mathbf{x})
=
-
\sum_i
\log
[
-f_i(\mathbf{x})
].
}
\]

Then solve

\[
\boxed{
\min_{\mathbf{x}}
t f_0(\mathbf{x})
+
\phi(\mathbf{x}),
}
\]

for increasing \(t\).

%%%%%%%%%%%%%%%%%%%%%%%%%%%%%%%%%%%%%%%%%%%%%%%%%%%%%%%%%%%%
\subsection{Barrier Central Path}
\label{subsec:convex-barrier-central-path}
%%%%%%%%%%%%%%%%%%%%%%%%%%%%%%%%%%%%%%%%%%%%%%%%%%%%%%%%%%%%

For each \(t>0\), let \(\mathbf{x}^*(t)\) solve the barrier problem.

The collection

\[
\boxed{
\{
\mathbf{x}^*(t):
t>0
\}
}
\]

is the central path.

As

\[
t\to\infty,
\]

the barrier solution approaches the constrained optimum under standard
conditions.

%%%%%%%%%%%%%%%%%%%%%%%%%%%%%%%%%%%%%%%%%%%%%%%%%%%%%%%%%%%%
\subsection{Barrier Multipliers}
\label{subsec:barrier-multipliers}
%%%%%%%%%%%%%%%%%%%%%%%%%%%%%%%%%%%%%%%%%%%%%%%%%%%%%%%%%%%%

For barrier solution \(\mathbf{x}^*(t)\), define

\[
\boxed{
\lambda_i
=
-
\frac{
1
}{
t f_i(\mathbf{x}^*(t))
}.
}
\]

Then

\[
\lambda_i>0
\]

and

\[
\boxed{
-\lambda_i f_i(\mathbf{x}^*)
=
\frac1t.
}
\]

Thus barrier methods satisfy an approximate complementary-slackness
condition.

%%%%%%%%%%%%%%%%%%%%%%%%%%%%%%%%%%%%%%%%%%%%%%%%%%%%%%%%%%%%
\subsection{Self-Concordance Preview}
\label{subsec:self-concordance-preview}
%%%%%%%%%%%%%%%%%%%%%%%%%%%%%%%%%%%%%%%%%%%%%%%%%%%%%%%%%%%%

Interior-point complexity theory relies on special barrier functions whose
third derivatives are controlled by their second derivatives.

Such functions are called self-concordant.

Self-concordance makes Newton steps predictable even near curved
boundaries.

This concept underlies polynomial-time complexity results for many convex
optimization classes.

%%%%%%%%%%%%%%%%%%%%%%%%%%%%%%%%%%%%%%%%%%%%%%%%%%%%%%%%%%%%
\subsection{Convex Relaxations}
\label{subsec:convex-relaxations}
%%%%%%%%%%%%%%%%%%%%%%%%%%%%%%%%%%%%%%%%%%%%%%%%%%%%%%%%%%%%

A relaxation replaces a difficult feasible set by a larger convex set.

If the original minimization problem is

\[
\boxed{
\min_{\mathbf{x}\in S}
f(\mathbf{x}),
}
\]

and

\[
S\subseteq C,
\]

then the relaxed problem is

\[
\boxed{
\min_{\mathbf{x}\in C}
f(\mathbf{x}).
}
\]

Its optimal value provides a lower bound on the original problem when
minimizing.

%%%%%%%%%%%%%%%%%%%%%%%%%%%%%%%%%%%%%%%%%%%%%%%%%%%%%%%%%%%%
\subsection{Why Convex Relaxations Matter}
\label{subsec:why-convex-relaxations}
%%%%%%%%%%%%%%%%%%%%%%%%%%%%%%%%%%%%%%%%%%%%%%%%%%%%%%%%%%%%

Convex relaxations are useful because they can:

\[
\boxed{
\text{produce rigorous bounds},
}
\]

\[
\boxed{
\text{provide good approximate solutions},
}
\]

and

\[
\boxed{
\text{support branch-and-bound methods}.
}
\]

They are especially important in integer and combinatorial optimization.

%%%%%%%%%%%%%%%%%%%%%%%%%%%%%%%%%%%%%%%%%%%%%%%%%%%%%%%%%%%%
\subsection{Quadratic Relaxation Example}
\label{subsec:quadratic-relaxation-example}
%%%%%%%%%%%%%%%%%%%%%%%%%%%%%%%%%%%%%%%%%%%%%%%%%%%%%%%%%%%%

Suppose a binary variable satisfies

\[
x\in\{0,1\}.
\]

A simple convex relaxation is

\[
\boxed{
0\leq x\leq1.
}
\]

The relaxed feasible set contains the original discrete set.

The relaxation may be much easier to optimize over.

%%%%%%%%%%%%%%%%%%%%%%%%%%%%%%%%%%%%%%%%%%%%%%%%%%%%%%%%%%%%
\subsection{Robust Optimization Preview}
\label{subsec:robust-convex-optimization-preview}
%%%%%%%%%%%%%%%%%%%%%%%%%%%%%%%%%%%%%%%%%%%%%%%%%%%%%%%%%%%%

Suppose a constraint contains uncertain data \(u\):

\[
f(\mathbf{x},u)\leq0.
\]

A robust constraint requires

\[
\boxed{
f(\mathbf{x},u)\leq0
\quad
\text{for every }
u\in\mathcal{U}.
}
\]

Under suitable uncertainty sets and functional structure, the robust
counterpart can remain convex.

%%%%%%%%%%%%%%%%%%%%%%%%%%%%%%%%%%%%%%%%%%%%%%%%%%%%%%%%%%%%
\subsection{Maximum Entropy Example}
\label{subsec:maximum-entropy-convex}
%%%%%%%%%%%%%%%%%%%%%%%%%%%%%%%%%%%%%%%%%%%%%%%%%%%%%%%%%%%%

For probabilities

\[
p_i\geq0,
\qquad
\sum_i p_i=1,
\]

entropy is

\[
\boxed{
H(\mathbf{p})
=
-
\sum_i
p_i\log p_i.
}
\]

Entropy is concave.

Therefore maximizing entropy subject to affine constraints is a convex
optimization problem in the standard sense of maximizing a concave
objective over a convex feasible set.

%%%%%%%%%%%%%%%%%%%%%%%%%%%%%%%%%%%%%%%%%%%%%%%%%%%%%%%%%%%%
\subsection{Portfolio Optimization}
\label{subsec:portfolio-convex-optimization}
%%%%%%%%%%%%%%%%%%%%%%%%%%%%%%%%%%%%%%%%%%%%%%%%%%%%%%%%%%%%

Let

\[
\mathbf{w}
\]

be portfolio weights and let

\[
\Sigma\succeq0
\]

be the covariance matrix.

Portfolio variance is

\[
\boxed{
\mathbf{w}^T
\Sigma
\mathbf{w}.
}
\]

A minimum-variance portfolio problem is

\[
\boxed{
\begin{aligned}
\min_{\mathbf{w}}
\quad&
\mathbf{w}^T
\Sigma
\mathbf{w}
\\
\text{subject to}
\quad&
\mathbf{1}^T\mathbf{w}=1,
\\
&
\boldsymbol{\mu}^T\mathbf{w}
\geq
r_0.
\end{aligned}
}
\]

This is a convex quadratic program.

%%%%%%%%%%%%%%%%%%%%%%%%%%%%%%%%%%%%%%%%%%%%%%%%%%%%%%%%%%%%
\subsection{Regularized Regression}
\label{subsec:regularized-regression-convex}
%%%%%%%%%%%%%%%%%%%%%%%%%%%%%%%%%%%%%%%%%%%%%%%%%%%%%%%%%%%%

Ridge regression solves

\[
\boxed{
\min_{\boldsymbol{\beta}}
\frac12
\|
X\boldsymbol{\beta}-\mathbf{y}
\|_2^2
+
\frac{\lambda}{2}
\|
\boldsymbol{\beta}
\|_2^2.
}
\]

This objective is convex.

Lasso solves

\[
\boxed{
\min_{\boldsymbol{\beta}}
\frac12
\|
X\boldsymbol{\beta}-\mathbf{y}
\|_2^2
+
\lambda
\|
\boldsymbol{\beta}
\|_1.
}
\]

This is convex but nonsmooth.

%%%%%%%%%%%%%%%%%%%%%%%%%%%%%%%%%%%%%%%%%%%%%%%%%%%%%%%%%%%%
\subsection{Logistic Regression as Convex Optimization}
\label{subsec:logistic-regression-convex}
%%%%%%%%%%%%%%%%%%%%%%%%%%%%%%%%%%%%%%%%%%%%%%%%%%%%%%%%%%%%

For binary responses \(y_i\in\{0,1\}\), logistic regression can be
estimated by minimizing the negative log-likelihood.

A common form is

\[
\boxed{
\sum_{i=1}^{n}
\left[
\log
\left(
1+e^{\mathbf{x}_i^T\boldsymbol{\beta}}
\right)
-
y_i
\mathbf{x}_i^T\boldsymbol{\beta}
\right].
}
\]

This objective is convex in \(\boldsymbol{\beta}\).

Regularization can also be added while preserving convexity for many
common penalties.

%%%%%%%%%%%%%%%%%%%%%%%%%%%%%%%%%%%%%%%%%%%%%%%%%%%%%%%%%%%%
\subsection{Support Vector Machines Preview}
\label{subsec:svm-convex-preview}
%%%%%%%%%%%%%%%%%%%%%%%%%%%%%%%%%%%%%%%%%%%%%%%%%%%%%%%%%%%%

A soft-margin support vector machine can be written as a convex quadratic
program:

\[
\boxed{
\begin{aligned}
\min_{\mathbf{w},b,\boldsymbol{\xi}}
\quad&
\frac12
\|\mathbf{w}\|_2^2
+
C
\sum_i
\xi_i
\\
\text{subject to}
\quad&
y_i
(
\mathbf{w}^T\mathbf{x}_i+b
)
\geq
1-\xi_i,
\\
&
\xi_i\geq0.
\end{aligned}
}
\]

Convexity guarantees global optimality.

%%%%%%%%%%%%%%%%%%%%%%%%%%%%%%%%%%%%%%%%%%%%%%%%%%%%%%%%%%%%
\subsection{Control Applications}
\label{subsec:convex-control-applications}
%%%%%%%%%%%%%%%%%%%%%%%%%%%%%%%%%%%%%%%%%%%%%%%%%%%%%%%%%%%%

Many control-design conditions can be written as linear matrix
inequalities.

For example, stability of

\[
x'=Ax
\]

can be certified by finding

\[
P\succ0
\]

such that

\[
\boxed{
A^TP+PA\prec0.
}
\]

This is a convex matrix inequality in \(P\).

Such formulations lead naturally to semidefinite programming.

%%%%%%%%%%%%%%%%%%%%%%%%%%%%%%%%%%%%%%%%%%%%%%%%%%%%%%%%%%%%
\subsection{Geometric Programming Preview}
\label{subsec:geometric-programming-preview}
%%%%%%%%%%%%%%%%%%%%%%%%%%%%%%%%%%%%%%%%%%%%%%%%%%%%%%%%%%%%

Some nonlinear problems involving positive variables and monomials or
posynomials can be transformed into convex problems using logarithmic
changes of variables.

These are called geometric programs.

For example, multiplicative structure becomes additive after taking
logarithms.

%%%%%%%%%%%%%%%%%%%%%%%%%%%%%%%%%%%%%%%%%%%%%%%%%%%%%%%%%%%%
\subsection{Convex Optimization Workflow}
\label{subsec:convex-optimization-workflow}
%%%%%%%%%%%%%%%%%%%%%%%%%%%%%%%%%%%%%%%%%%%%%%%%%%%%%%%%%%%%

A practical workflow is:

\begin{enumerate}
    \item identify the decision variables;
    \item write the objective and constraints;
    \item determine the domain;
    \item verify that the feasible set is convex;
    \item verify that the objective is convex for minimization;
    \item verify that inequality functions are convex;
    \item verify that equality constraints are affine;
    \item use algebraic convexity-preservation rules when possible;
    \item check strict or strong convexity if uniqueness matters;
    \item derive KKT conditions when useful;
    \item check Slater's condition for strong duality;
    \item formulate the dual problem when it provides insight;
    \item identify special structure such as LP, QP, SOCP, or SDP;
    \item choose an appropriate first-order or interior-point method;
    \item scale the problem;
    \item verify feasibility and duality gap;
    \item interpret multipliers and sensitivity.
\end{enumerate}

%%%%%%%%%%%%%%%%%%%%%%%%%%%%%%%%%%%%%%%%%%%%%%%%%%%%%%%%%%%%
\subsection{Recognizing Convex Structure}
\label{subsec:recognizing-convex-structure}
%%%%%%%%%%%%%%%%%%%%%%%%%%%%%%%%%%%%%%%%%%%%%%%%%%%%%%%%%%%%

A problem may look nonlinear but still be convex.

Useful questions include:

\begin{enumerate}
    \item Is the objective a known convex function?
    \item Is it a nonnegative sum of convex functions?
    \item Is it a convex function composed with an affine mapping?
    \item Are inequality functions convex?
    \item Are equality constraints affine?
    \item Can an epigraph transformation reveal convex structure?
    \item Can a norm or cone formulation simplify the model?
\end{enumerate}

Recognizing convexity is often more valuable than manipulating the model
into an arbitrary nonlinear-programming form.

%%%%%%%%%%%%%%%%%%%%%%%%%%%%%%%%%%%%%%%%%%%%%%%%%%%%%%%%%%%%
\subsection{Common Errors}
\label{subsec:convex-optimization-common-errors}
%%%%%%%%%%%%%%%%%%%%%%%%%%%%%%%%%%%%%%%%%%%%%%%%%%%%%%%%%%%%

\subsubsection{Calling Every Nonlinear Problem Nonconvex}

Nonlinearity does not imply nonconvexity.

For example,

\[
x^2
\]

is nonlinear but convex.

%%%%%%%%%%%%%%%%%%%%%%%%%%%%%%%%%%%%%%%%%%%%%%%%%%%%%%%%%%%%
\subsubsection{Checking Only the Objective}

For a convex optimization problem, the feasible region must also be
convex.

A convex objective over a nonconvex set is not a convex optimization
problem.

%%%%%%%%%%%%%%%%%%%%%%%%%%%%%%%%%%%%%%%%%%%%%%%%%%%%%%%%%%%%
\subsubsection{Allowing Nonaffine Equality Constraints}

A constraint such as

\[
x^2+y^2=1
\]

does not define a convex feasible set.

Standard convex optimization permits affine equality constraints.

%%%%%%%%%%%%%%%%%%%%%%%%%%%%%%%%%%%%%%%%%%%%%%%%%%%%%%%%%%%%
\subsubsection{Assuming Positive Semidefinite Hessian at One Point Proves Convexity}

To prove convexity by the Hessian criterion, one needs

\[
\boxed{
\nabla^2f(\mathbf{x})
\succeq0
}
\]

throughout the convex domain.

Checking one point is insufficient.

%%%%%%%%%%%%%%%%%%%%%%%%%%%%%%%%%%%%%%%%%%%%%%%%%%%%%%%%%%%%
\subsubsection{Confusing Strict and Strong Convexity}

Strong convexity implies strict convexity.

Strict convexity does not necessarily imply strong convexity.

%%%%%%%%%%%%%%%%%%%%%%%%%%%%%%%%%%%%%%%%%%%%%%%%%%%%%%%%%%%%
\subsubsection{Assuming Convex Functions Are Always Differentiable}

Functions such as

\[
|x|
\]

and

\[
\|\mathbf{x}\|_1
\]

are convex but nonsmooth.

Subgradients are required at nondifferentiable points.

%%%%%%%%%%%%%%%%%%%%%%%%%%%%%%%%%%%%%%%%%%%%%%%%%%%%%%%%%%%%
\subsubsection{Forgetting Domain Restrictions}

The function

\[
-\log x
\]

is convex only on its domain

\[
\boxed{
x>0.
}
\]

Convexity statements must include the domain.

%%%%%%%%%%%%%%%%%%%%%%%%%%%%%%%%%%%%%%%%%%%%%%%%%%%%%%%%%%%%
\subsubsection{Using Composition Rules Without Checking Monotonicity}

A convex outer function composed with a convex inner function is not
automatically convex.

Monotonicity and domain conditions matter.

%%%%%%%%%%%%%%%%%%%%%%%%%%%%%%%%%%%%%%%%%%%%%%%%%%%%%%%%%%%%
\subsubsection{Assuming Slater's Condition Is Necessary for Strong Duality}

Slater's condition is an important sufficient condition.

Strong duality may still hold even when Slater's condition fails.

%%%%%%%%%%%%%%%%%%%%%%%%%%%%%%%%%%%%%%%%%%%%%%%%%%%%%%%%%%%%
\subsubsection{Confusing Weak and Strong Duality}

Weak duality always gives

\[
\boxed{
d^*\leq p^*.
}
\]

Strong duality additionally states

\[
\boxed{
d^*=p^*.
}
\]

%%%%%%%%%%%%%%%%%%%%%%%%%%%%%%%%%%%%%%%%%%%%%%%%%%%%%%%%%%%%
\subsubsection{Ignoring Complementary Slackness}

At a convex KKT point,

\[
\boxed{
\lambda_i f_i(x^*)=0.
}
\]

A strictly inactive inequality must have zero multiplier.

%%%%%%%%%%%%%%%%%%%%%%%%%%%%%%%%%%%%%%%%%%%%%%%%%%%%%%%%%%%%
\subsubsection{Assuming Any KKT Point Is Globally Optimal}

Global sufficiency requires convex problem structure.

For nonconvex problems, KKT conditions alone do not give global
optimality.

%%%%%%%%%%%%%%%%%%%%%%%%%%%%%%%%%%%%%%%%%%%%%%%%%%%%%%%%%%%%
\subsubsection{Treating a Convex Relaxation as the Original Problem}

A relaxation enlarges the feasible set.

Its solution may be infeasible for the original problem.

Its objective value is primarily a bound until a feasible original
solution is recovered.

%%%%%%%%%%%%%%%%%%%%%%%%%%%%%%%%%%%%%%%%%%%%%%%%%%%%%%%%%%%%
\subsubsection{Ignoring Conditioning}

A convex problem may still be numerically difficult.

Convexity guarantees global structure, not necessarily good numerical
conditioning.

%%%%%%%%%%%%%%%%%%%%%%%%%%%%%%%%%%%%%%%%%%%%%%%%%%%%%%%%%%%%
\subsection{Applications}
\label{subsec:convex-optimization-applications}
%%%%%%%%%%%%%%%%%%%%%%%%%%%%%%%%%%%%%%%%%%%%%%%%%%%%%%%%%%%%

\begin{applicationbox}

\textbf{Statistics.}
Least squares, logistic regression, regularized regression, covariance
estimation, and many likelihood problems are convex.

\medskip

\textbf{Machine Learning.}
Support vector machines, sparse regression, and many classification and
estimation problems rely on convex optimization.

\medskip

\textbf{Signal Processing.}
Denoising, sparse reconstruction, compressed sensing, and inverse
problems frequently use convex norms and regularization.

\medskip

\textbf{Finance.}
Mean-variance portfolio optimization and many risk-constrained allocation
problems are convex quadratic or conic programs.

\medskip

\textbf{Control Theory.}
Linear matrix inequalities and Lyapunov stability conditions lead to
semidefinite programs.

\medskip

\textbf{Engineering Design.}
Norm constraints, quadratic models, robust design, and geometric programs
often admit convex formulations.

\medskip

\textbf{Operations Research.}
Linear programming, network models, convex resource allocation, and
relaxations of discrete problems all fit within convex optimization.

\medskip

\textbf{Inverse Problems.}
Convex regularizers transform ill-posed estimation into stable
optimization problems.

\end{applicationbox}

%%%%%%%%%%%%%%%%%%%%%%%%%%%%%%%%%%%%%%%%%%%%%%%%%%%%%%%%%%%%
\subsection{Exercises}
\label{subsec:convex-optimization-exercises}
%%%%%%%%%%%%%%%%%%%%%%%%%%%%%%%%%%%%%%%%%%%%%%%%%%%%%%%%%%%%

\begin{exercise}[1]
Define a convex set.
\end{exercise}

\begin{exercise}[1]
Define a convex combination.
\end{exercise}

\begin{exercise}[1]
Define the convex hull.
\end{exercise}

\begin{exercise}[1]
Define a convex function.
\end{exercise}

\begin{exercise}[1]
Define strict convexity.
\end{exercise}

\begin{exercise}[1]
Define strong convexity.
\end{exercise}

\begin{exercise}[1]
Define the epigraph of a function.
\end{exercise}

\begin{exercise}[1]
Define a subgradient.
\end{exercise}

\begin{exercise}[1]
Define a convex cone.
\end{exercise}

\begin{exercise}[1]
Define the Lagrange dual function.
\end{exercise}

\begin{exercise}[1]
State weak duality.
\end{exercise}

\begin{exercise}[1]
State Slater's condition.
\end{exercise}

\begin{exercise}[2]
Determine whether

\[
\{
(x,y):
x+y\leq4
\}
\]

is convex.
\end{exercise}

\begin{exercise}[2]
Determine whether

\[
\{
(x,y):
x^2+y^2\leq1
\}
\]

is convex.
\end{exercise}

\begin{exercise}[2]
Determine whether

\[
\{
(x,y):
x^2+y^2=1
\}
\]

is convex.
\end{exercise}

\begin{exercise}[2]
Determine whether

\[
f(x)=x^4
\]

is convex.
\end{exercise}

\begin{exercise}[2]
Determine whether

\[
f(x)=e^x
\]

is convex.
\end{exercise}

\begin{exercise}[2]
Determine whether

\[
f(x)=-\log x
\]

is convex on \(x>0\).
\end{exercise}

\begin{exercise}[2]
Compute the Hessian of

\[
f(x,y)
=
x^2+xy+y^2.
\]
\end{exercise}

\begin{exercise}[2]
Determine whether the previous function is convex.
\end{exercise}

\begin{exercise}[2]
Find the subdifferential of

\[
|x|
\]

at \(x=0\).
\end{exercise}

\begin{exercise}[2]
Determine whether

\[
0\in\partial|x|
\]

at \(x=0\), and interpret the result.
\end{exercise}

\begin{exercise}[3]
Prove that every affine set is convex.
\end{exercise}

\begin{exercise}[3]
Prove that intersections of convex sets are convex.
\end{exercise}

\begin{exercise}[3]
Prove that every norm ball is convex.
\end{exercise}

\begin{exercise}[3]
Prove that a nonnegative weighted sum of convex functions is convex.
\end{exercise}

\begin{exercise}[3]
Prove that affine composition preserves convexity.
\end{exercise}

\begin{exercise}[3]
Prove that the pointwise maximum of convex functions is convex.
\end{exercise}

\begin{exercise}[3]
Prove that the epigraph of a convex function is convex.
\end{exercise}

\begin{exercise}[3]
Prove that every local minimizer of a convex function on a convex set is
global.
\end{exercise}

\begin{exercise}[3]
Prove that a strictly convex function has at most one minimizer.
\end{exercise}

\begin{exercise}[3]
Use the Hessian to test convexity of

\[
f(x,y)
=
3x^2+2xy+2y^2.
\]
\end{exercise}

\begin{exercise}[3]
Write

\[
\min_x |x|
\]

in epigraph form.
\end{exercise}

\begin{exercise}[3]
Write

\[
\min_x
\max
\{
2x+1,
-x+3
\}
\]

in epigraph form.
\end{exercise}

\begin{exercise}[3]
Show that the second-order cone is convex.
\end{exercise}

\begin{exercise}[3]
Show that the positive-semidefinite cone is convex.
\end{exercise}

\begin{exercise}[3]
Formulate a Euclidean norm minimization problem as an SOCP.
\end{exercise}

\begin{exercise}[3]
Formulate a convex quadratic program with three variables and two affine
constraints.
\end{exercise}

\begin{exercise}[3]
Derive the dual function for

\[
\min_x x^2
\]

subject to

\[
1-x\leq0.
\]
\end{exercise}

\begin{exercise}[3]
Solve the previous problem using KKT conditions.
\end{exercise}

\begin{exercise}[3]
Verify strong duality for the previous problem.
\end{exercise}

\begin{exercise}[3]
Check Slater's condition for

\[
x^2-4<0.
\]
\end{exercise}

\begin{exercise}[3]
Find the projection of \((3,4)\) onto the unit Euclidean ball.
\end{exercise}

\begin{exercise}[3]
Find the projection of a point onto a half-space.
\end{exercise}

\begin{exercise}[3]
Compute the proximal operator of

\[
\lambda|x|
\]

and derive soft thresholding.
\end{exercise}

\begin{exercise}[3]
Derive the projected-gradient update for a box-constrained problem.
\end{exercise}

\begin{exercise}[3]
Formulate ridge regression as a convex optimization problem.
\end{exercise}

\begin{exercise}[3]
Formulate lasso regression as a convex optimization problem.
\end{exercise}

\begin{exercise}[3]
Show that the logistic-regression negative log-likelihood is convex.
\end{exercise}

\begin{exercise}[3]
Formulate a minimum-variance portfolio problem.
\end{exercise}

\begin{exercise}[3]
Explain why the Lyapunov inequality

\[
A^TP+PA\prec0
\]

is affine in \(P\).
\end{exercise}

\begin{exercise}[4]
Prove the first-order characterization of convexity.
\end{exercise}

\begin{exercise}[4]
Prove the second-order Hessian characterization of convexity.
\end{exercise}

\begin{exercise}[4]
Study quasiconvex functions and quasiconvex optimization.
\end{exercise}

\begin{exercise}[4]
Study perspective transformations of convex functions.
\end{exercise}

\begin{exercise}[4]
Study supporting-hyperplane theorems.
\end{exercise}

\begin{exercise}[4]
Study separating-hyperplane theorems.
\end{exercise}

\begin{exercise}[4]
Prove the subgradient optimality condition.
\end{exercise}

\begin{exercise}[4]
Study normal cones and constrained convex optimality.
\end{exercise}

\begin{exercise}[4]
Derive Lagrange duality from lower-bounding arguments.
\end{exercise}

\begin{exercise}[4]
Prove weak duality for general convex programs.
\end{exercise}

\begin{exercise}[4]
Study strong duality under Slater's condition.
\end{exercise}

\begin{exercise}[4]
Derive KKT necessity and sufficiency for convex problems.
\end{exercise}

\begin{exercise}[4]
Study convex conjugates.
\end{exercise}

\begin{exercise}[4]
Prove the Fenchel--Young inequality.
\end{exercise}

\begin{exercise}[4]
Study Fenchel duality.
\end{exercise}

\begin{exercise}[4]
Study proximal operators and monotone mappings.
\end{exercise}

\begin{exercise}[4]
Derive convergence rates for smooth convex gradient descent.
\end{exercise}

\begin{exercise}[4]
Derive linear convergence under strong convexity.
\end{exercise}

\begin{exercise}[4]
Study Nesterov acceleration.
\end{exercise}

\begin{exercise}[4]
Study proximal-gradient convergence.
\end{exercise}

\begin{exercise}[4]
Study primal-dual first-order algorithms.
\end{exercise}

\begin{exercise}[4]
Study conic duality.
\end{exercise}

\begin{exercise}[4]
Derive the dual of an SOCP.
\end{exercise}

\begin{exercise}[4]
Derive the dual of a semidefinite program.
\end{exercise}

\begin{exercise}[4]
Study self-concordant barriers.
\end{exercise}

\begin{exercise}[4]
Study polynomial-time complexity of interior-point methods.
\end{exercise}

\begin{exercise}[4]
Construct convex relaxations for binary quadratic problems.
\end{exercise}

\begin{exercise}[4]
Study semidefinite relaxations of combinatorial optimization problems.
\end{exercise}

%%%%%%%%%%%%%%%%%%%%%%%%%%%%%%%%%%%%%%%%%%%%%%%%%%%%%%%%%%%%
\subsection{Conceptual Summary}
\label{subsec:convex-optimization-summary}
%%%%%%%%%%%%%%%%%%%%%%%%%%%%%%%%%%%%%%%%%%%%%%%%%%%%%%%%%%%%

A convex optimization problem has form

\[
\boxed{
\begin{aligned}
\min_{\mathbf{x}}
\quad&
f_0(\mathbf{x})
\\
\text{subject to}
\quad&
f_i(\mathbf{x})\leq0,
\\
&
A\mathbf{x}=\mathbf{b},
\end{aligned}
}
\]

where the objective and inequality functions are convex.

A set \(C\) is convex when

\[
\boxed{
\theta\mathbf{x}
+
(1-\theta)\mathbf{y}
\in C
}
\]

for all

\[
\mathbf{x},\mathbf{y}\in C
\]

and

\[
0\leq\theta\leq1.
\]

A differentiable function is convex when

\[
\boxed{
f(\mathbf{y})
\geq
f(\mathbf{x})
+
\nabla f(\mathbf{x})^T
(
\mathbf{y}-\mathbf{x}
).
}
\]

For twice-differentiable functions,

\[
\boxed{
\nabla^2f(\mathbf{x})\succeq0
}
\]

throughout the domain implies convexity.

For convex problems,

\[
\boxed{
\text{local optimum}
=
\text{global optimum}.
}
\]

The Lagrange dual function is

\[
\boxed{
q(\lambda,\nu)
=
\inf_x
\mathcal{L}(x,\lambda,\nu).
}
\]

Weak duality gives

\[
\boxed{
d^*\leq p^*.
}
\]

Under suitable conditions such as Slater's condition,

\[
\boxed{
d^*=p^*.
}
\]

Convex optimization therefore combines

\[
\boxed{
\text{geometry}
+
\text{duality}
+
\text{linear algebra}
+
\text{analysis}
+
\text{efficient algorithms}.
}
\]

%%%%%%%%%%%%%%%%%%%%%%%%%%%%%%%%%%%%%%%%%%%%%%%%%%%%%%%%%%%%
\subsection{Section Connections}
\label{subsec:convex-optimization-connections}
%%%%%%%%%%%%%%%%%%%%%%%%%%%%%%%%%%%%%%%%%%%%%%%%%%%%%%%%%%%%

\chapterconnection{
Convex Optimization identifies the globally tractable structure hidden
inside many nonlinear models. Convex sets, subgradients, duality,
conic formulations, proximal methods, and interior-point methods connect
the general nonlinear theory of Section~\ref{sec:nonlinear-programming}
with linear programming and modern large-scale optimization. Convex
relaxations also provide essential lower bounds for discrete problems.
The next section, Integer Programming, introduces decision variables that
must take integral values and develops branch-and-bound, cutting planes,
relaxations, and mixed-integer modeling.
}

%%%%%%%%%%%%%%%%%%%%%%%%%%%%%%%%%%%%%%%%%%%%%%%%%%%%%%%%%%%%
% SECTION SOURCE TRACKING
%%%%%%%%%%%%%%%%%%%%%%%%%%%%%%%%%%%%%%%%%%%%%%%%%%%%%%%%%%%%

%------------------------------------------------------------
% 10.4 Convex Optimization
%------------------------------------------------------------
%
% Source basis:
%   - Standard convex analysis.
%   - Standard convex optimization and duality theory.
%   - Standard conic and first-order optimization methods.
%
% Used for:
%   - convex combinations;
%   - convex sets;
%   - affine sets and norm balls;
%   - convex hulls;
%   - intersections of convex sets;
%   - convex, strict, and strongly convex functions;
%   - first- and second-order convexity criteria;
%   - epigraphs and sublevel sets;
%   - convexity-preservation rules;
%   - affine composition;
%   - pointwise maxima;
%   - perspective functions;
%   - convex cones and dual cones;
%   - second-order cones;
%   - SOCP;
%   - positive-semidefinite cones;
%   - semidefinite programming;
%   - convex quadratic programming;
%   - projections onto convex sets;
%   - supporting and separating hyperplanes;
%   - subgradients and subdifferentials;
%   - normal cones;
%   - constrained convex optimality;
%   - Lagrange dual functions;
%   - weak and strong duality;
%   - Slater's condition;
%   - KKT sufficiency;
%   - convex conjugates;
%   - Fenchel--Young inequality;
%   - proximal operators;
%   - soft thresholding;
%   - gradient, subgradient, projected-gradient, and proximal-gradient
%     methods;
%   - accelerated and coordinate methods preview;
%   - interior-point methods;
%   - central paths;
%   - self-concordance preview;
%   - convex relaxations;
%   - robust optimization preview;
%   - statistical, machine-learning, financial, signal-processing, and
%     control applications;
%   - exercises.
%
% Notes:
%   - Integer Programming is developed next in Section 10.5.
%   - Combinatorial Optimization follows in Section 10.6.
%   - Numerical optimization methods are developed more deeply in
%     Chapter 11.
%
%%%%%%%%%%%%%%%%%%%%%%%%%%%%%%%%%%%%%%%%%%%%%%%%%%%%%%%%%%%%